\chapter{Có và Không}
\markboth{Có và Không}{Having and Not Having}

\section{Một người có nhà đẹp nhưng hiếm khi ở nhà.}
\begin{truyen}{Một người có nhà đẹp nhưng hiếm khi ở nhà.}{Having a house is not always having a home.}
\chuhoa{A}{nh làm việc rất nhiều để mua được căn nhà rộng như mơ. Nhưng khi có rồi, anh lại đi sớm về khuya và chỉ xem nó như chỗ ngủ. Cái có trong giấy tờ chưa chắc là cái có trong đời sống thật.}
\textit{(He worked very hard to buy the spacious house he once dreamed of. But after getting it, he left early and returned late, using it mostly as a place to sleep. What you own on paper is not always what you truly have in life.)}
\end{truyen}
\begin{baihoc}
Có một thứ trong tay chưa chắc đã có nó trong lòng.
\textit{(Holding something in your hand does not mean holding it in your life.)}
\end{baihoc}
\begin{ghinhoanh}
\textbf{Short English to remember:}

Owning a house is not always having a home.
\end{ghinhoanh}
\begin{tuhocnhanh}
1. Em đang có điều gì mà thật ra chưa sống được cùng nó? \textit{(What do you own but have not really lived with yet?)}

2. Điều gì làm một thứ trở thành của mình thật sự? \textit{(What makes something truly yours?)}
\end{tuhocnhanh}
\ngancach

\section{Một học sinh không có điện thoại đắt tiền nhưng có giờ tự học rất đều.}
\begin{truyen}{Một học sinh không có điện thoại đắt tiền nhưng có giờ tự học rất đều.}{Not having much can protect focus.}
\chuhoa{B}{ạn bè thấy em thiệt thòi vì không có đồ mới như họ. Nhưng chính sự thiếu ấy khiến em ít xao động hơn, giữ được nhịp học bền hơn. Không có một thứ đôi khi lại giữ được điều quan trọng hơn.}
\textit{(Friends thought he was disadvantaged because he lacked the newest devices. Yet that very lack kept him less distracted and more consistent in study. Not having one thing can sometimes protect something more important.)}
\end{truyen}
\begin{baihoc}
Thiếu vật chất không luôn là nghèo; có khi nó giữ cho lòng mình bớt phân tán.
\textit{(Material lack is not always poverty; sometimes it helps the mind stay less scattered.)}
\end{baihoc}
\begin{ghinhoanh}
\textbf{Short English to remember:}

Less stuff can mean more focus.
\end{ghinhoanh}
\begin{tuhocnhanh}
1. Có thứ gì em không có nhưng lại giúp em đỡ phân tâm hơn? \textit{(Is there anything you lack that actually helps you stay focused?)}

2. Em đang xem thiếu thốn nào quá bi quan không? \textit{(Are you viewing any lack too pessimistically?)}
\end{tuhocnhanh}
\ngancach

\section{Một người có rất nhiều mối quen nhưng không có ai để gọi lúc buồn.}
\begin{truyen}{Một người có rất nhiều mối quen nhưng không có ai để gọi lúc buồn.}{Having many contacts is not having connection.}
\chuhoa{N}{goài mặt anh rất rộng quan hệ, đi đâu cũng có người chào hỏi. Nhưng đến lúc thật sự xuống tinh thần, anh mở danh bạ rất lâu mà không biết gọi cho ai. Có nhiều người quanh mình không đồng nghĩa với có người ở cùng mình.}
\textit{(Outwardly, he knew many people and was greeted everywhere. But when he truly felt low, he stared at his contact list and did not know whom to call. Having many people around you does not mean having someone with you.)}
\end{truyen}
\begin{baihoc}
Đông chưa chắc là đầy; có khi rất đông mà vẫn rất trống.
\textit{(Crowded is not always full; sometimes it is crowded and still empty.)}
\end{baihoc}
\begin{ghinhoanh}
\textbf{Short English to remember:}

Many contacts are not the same as connection.
\end{ghinhoanh}
\begin{tuhocnhanh}
1. Em đang xây danh sách quen biết hay đang nuôi những người thật sự gần? \textit{(Are you building contacts or nurturing real closeness?)}

2. Ai là người em có thể gọi lúc không ổn? \textit{(Who can you call when you are not okay?)}
\end{tuhocnhanh}
\ngancach

\section{Một cô gái không có ngoại hình nổi bật nhưng có sự dịu dàng hiếm.}
\begin{truyen}{Một cô gái không có ngoại hình nổi bật nhưng có sự dịu dàng hiếm.}{What you lack outside may reveal another gift.}
\chuhoa{C}{ô từng tự ti vì mình không rực rỡ như người khác. Nhưng chính sự chân thành, nhẹ nhàng và biết lắng nghe của cô lại khiến nhiều người nhớ lâu hơn cả vẻ ngoài. Không có một thứ không có nghĩa là không có giá trị.}
\textit{(She used to feel insecure because she was not striking in appearance. Yet her sincerity, gentleness, and listening made people remember her more deeply than looks. Lacking one thing does not mean lacking value.)}
\end{truyen}
\begin{baihoc}
Con người thường nhìn cái mình không có quá kỹ mà quên cái quý mình đang có.
\textit{(People often stare too long at what they lack and forget what is precious in them already.)}
\end{baihoc}
\begin{ghinhoanh}
\textbf{Short English to remember:}

Lacking one thing does not mean lacking worth.
\end{ghinhoanh}
\begin{tuhocnhanh}
1. Em đang thiếu điều gì mà vì thế quên mất điểm mạnh của mình? \textit{(What do you feel you lack so strongly that you forget your strengths?)}

2. Điều quý nào trong em đáng được nhìn lại? \textit{(What precious quality in you deserves to be seen again?)}
\end{tuhocnhanh}
\ngancach

\section{Một người có mọi tiện nghi nhưng không có thời gian ngồi yên.}
\begin{truyen}{Một người có mọi tiện nghi nhưng không có thời gian ngồi yên.}{Comfort without stillness is still a kind of lack.}
\chuhoa{M}{ọi thứ quanh anh đều đủ: xe tốt, nhà đẹp, thiết bị mới. Nhưng anh gần như không có một buổi chiều nào thảnh thơi. Càng đủ tiện nghi, anh càng thấy mình thiếu một điều rất nền tảng: khoảng trống để sống chậm lại.}
\textit{(Everything around him was sufficient: a good car, nice home, new devices. But he had almost no quiet afternoon. The more comfort he had, the more he felt the lack of something basic: space to slow down.)}
\end{truyen}
\begin{baihoc}
Có nhiều tiện nghi mà không có khoảng thở vẫn là một kiểu thiếu.
\textit{(Having many comforts without breathing space is still a kind of lack.)}
\end{baihoc}
\begin{ghinhoanh}
\textbf{Short English to remember:}

Comfort without stillness is still lack.
\end{ghinhoanh}
\begin{tuhocnhanh}
1. Điều gì trong đời em đang đủ mà vẫn chưa thật sự đủ? \textit{(What in your life is abundant and still somehow not enough?)}

2. Em cần thêm tiện nghi hay cần thêm khoảng thở? \textit{(Do you need more comfort or more breathing room?)}
\end{tuhocnhanh}
\ngancach